为了展示像~$\Th{Q}$ 这样的理论如何能``谈论''可计算函数——尤其是通过哥德尔数编码的可证性——我们确立了 $\Th{Q}$ \textbf{表示}所有可计算函数。所谓``$\Th{Q}$ 表示 $f(n)$''，是指存在 $\Lang{L_A}$ 中的一个公式~$!A_f(x, y)$ 表达 $f(x) = y$，并且 $\Th{Q}$ 能够证明这一点。这进而意味着，只要 $f(n) = m$，就有 $\Th{Q} \Proves !A_f(\num{n},
\num{m})$ 和 $\Th{Q} \Proves \lforall[y][(!A_f(\num{n}, y) \lif y =
\num{m})]$。（这里，$\num{n}$ 是~$n$ 的\textbf{标准数码}，即项 $\Obj 0^{\prime\dots\prime}$，带有 $n$ 个~$\prime$。项~$\num{n}$ 在标准模型~$\Struct{N}$ 中指称数字~$n$，因此它是在 $\Lang{L_A}$ 中表示数字~$n$ 的一种便捷方式。）为了证明 $\Th{Q}$ 表示所有可计算函数，我们回到对可计算函数的刻画：它们可以从 $\Zero$、$\Succ$ 以及投影函数出发，通过复合、原始递归和正则极小化来定义。虽然证明基本函数可表示、以及由可表示函数经复合和正则极小化定义的函数也可表示，都相对容易，但原始递归则更为困难。我们证明了，如果允许少量额外的基本函数（即加法、乘法以及等词~$=$ 的特征函数），实际上就可以避免以原始递归方式定义函数。这就需要用到\textbf{beta 函数}，它使我们能够以初步的方式处理数的序列，并且它本身可以在不使用原始递归的情况下加以定义。